\section{Discusión}
\label{sec:discusion}

\paragraph{¿Es creíble una exactitud del 99\,\%?}
Una exactitud tan alta debe tratarse con cautela, porque la causa más habitual de resultados
demasiado buenos es una fuga de información entre entrenamiento y test \citep{kaufman2012leakage}.
Nuestro plan de trabajo consideraba de hecho un resultado por encima del 98\,\% como una señal de
alarma. Hemos descartado las fugas más probables:
\begin{enumerate*}[label=(\roman*)]
  \item la partición es por persona y una prueba automática confirma que no hay personas
    compartidas;
  \item las estadísticas de normalización se calculan solo con entrenamiento;
  \item todas las decisiones se tomaron con validación y el test se evaluó una sola vez por
    modelo (con la salvedad, explicada en la \cref{sec:protocolo}, de que el modelo final se
    cambió tras la búsqueda de hiperparámetros, cuando el test ya se había evaluado con el
    anterior); y
  \item validación y test, con personas distintas, dan resultados casi idénticos.
\end{enumerate*}
Creemos que la exactitud alta se explica más bien por la naturaleza de la tarea tal como la
hemos planteado. Los recortes se obtienen de las cajas anotadas a mano, así que la mano siempre
está centrada y completa. Los gestos de HaGRID están posados deliberadamente para la cámara, y
la mayoría de las clases tienen formas claramente distintas. La dificultad residual se concentra en
\cls{stop}/\cls{palm} y \cls{like}/\cls{call} y en las etiquetas ruidosas de \cls{no\_gesture}.

Esta interpretación delimita también el alcance del resultado. La cifra del 99\,\% estima el
rendimiento del \emph{clasificador} sobre recortes de calidad HaGRID con cajas perfectas, no el
del sistema completo con una webcam. En la demo, las cajas vendrán de MediaPipe y estarán peor
centradas, la cámara y la iluminación serán otras y aparecerán posturas de transición entre
gestos que HaGRID no contiene. Es esperable una caída de rendimiento por este cambio de dominio
que nuestro experimento no cuantifica.

\paragraph{Qué aporta cada componente.}
La ablación muestra un resultado menos uniforme de lo que se suele esperar: de los tres
componentes, solo la normalización por lotes tiene un efecto claro sobre la exactitud final. Su
papel es sobre todo de \emph{optimización}: sin BN la red aprende mucho más despacio y no llega a
converger en 50 épocas, en línea con la motivación original de BN como técnica para estabilizar y
acelerar el entrenamiento \citep{ioffe2015bn}. El \emph{dropout} y el aumento de datos son, en
cambio, técnicas de \emph{regularización}, y su efecto se ve en la dinámica del entrenamiento (más
distancia entre entrenamiento y validación y parada más temprana sin ellas) más que en la
exactitud final. Nuestra interpretación es que, en este problema, el modelo completo sobreajusta
poco: tiene muchos ejemplos por clase, clases equilibradas y la parada anticipada ya lo protege.
No es un argumento para eliminar el \emph{dropout} o el aumento de datos. Ambos deberían ganar
importancia ante el cambio de dominio a la webcam (otras cámaras, iluminación y encuadres), que
nuestro conjunto de test no representa; el aumento de datos se diseñó precisamente para simular
esas variaciones.

\paragraph{Sensibilidad a los hiperparámetros.}
La búsqueda muestra que, alrededor de la configuración inicial, el resultado depende poco de los
hiperparámetros: nueve cambios de tasa de aprendizaje, lote, capacidad y activación mueven la
exactitud de validación menos de medio punto, y casi siempre dentro de la variación entre semillas.
El único cambio elegido, el lote de 32, mejoró la validación pero no el test, lo que ilustra el
riesgo de elegir entre muchas configuraciones con un único conjunto de validación pequeño. El
resultado más útil de la búsqueda es otro: con cuatro veces menos parámetros la red rinde igual,
así que la capacidad de GestureCNN no es lo que limita la exactitud.

\paragraph{CNN propia frente a modelo preentrenado.}
La comparación con ResNet18 matiza la idea de que un modelo más grande y preentrenado es siempre
mejor. Con \NTrain{} imágenes de entrenamiento, una tarea bien acotada y entradas de
$128\times128$, una red de \ParamsFinal{} parámetros entrenada desde cero queda a una distancia
que no podemos distinguir del ruido estadístico. El preentrenamiento sí acelera mucho la
convergencia, y podría marcar más diferencia con menos datos o en el cambio de dominio a la
webcam, algo que no hemos medido. Por otra parte, el análisis de coste muestra que el número de
parámetros es una mala medida del coste de inferencia: la red pequeña es la que más operaciones
realiza, porque trabaja a mayor resolución en sus primeras capas.

\paragraph{Limitaciones.}
\begin{itemize}
  \item \textbf{Cajas anotadas.} Todo el entrenamiento y la evaluación usan las cajas de HaGRID.
    No hemos medido el rendimiento con las cajas de un detector ni con vídeo propio, que es el
    escenario real de uso.
  \item \textbf{Resolución.} Al trabajar con la versión de 384\,px, la mayoría de las manos se
    amplían para llegar a $128\times128$ (mediana del lado de la caja: \LadoCajaMediana), lo que
    elimina detalle fino como la posición exacta del meñique, relevante en \cls{like}/\cls{call}.
  \item \textbf{Pocas repeticiones.} Dos semillas por variante bastan para ver la magnitud de la
    variación entre ejecuciones, pero no para estimar con precisión diferencias de pocas décimas.
    La ResNet18 se entrenó una sola vez.
  \item \textbf{Búsqueda de hiperparámetros limitada.} La búsqueda cambia un hiperparámetro cada
    vez, con una sola semilla por configuración, de modo que no explora interacciones y sus
    diferencias son del orden del ruido entre ejecuciones. Además, se hizo solo sobre el modelo
    completo: es posible que otra tasa de aprendizaje mejore alguna variante de la ablación, en
    particular la red sin BN.
  \item \textbf{Ruido de etiqueta en la clase nula.} Algunos \cls{no\_gesture} son en realidad
    gestos (\cref{sec:gradcam}), lo que penaliza al modelo por predicciones correctas y hace que
    la exactitud medida sea una cota algo pesimista en esa clase.
  \item \textbf{Una sola fuente de datos.} Las personas, los fondos y la forma de grabar son los
    de HaGRID; no sabemos cómo se comporta el modelo con otras poblaciones o cámaras.
\end{itemize}
